\section*{अध्याय \ref{ch:b3:renal-osmoregulation} --- वृक्क-शरीरक्रिया और परासरण-नियमन}

\begin{solution}{exo:b3:renal-osmoregulation:1}
\omterm{def:b3:renal-osmoregulation:nephron}{केशिकागुच्छ}: प्रोटीन छोड़कर प्लाज़्मा छानता है। \omterm{def:b3:renal-osmoregulation:nephron}{समीपस्थ नलिका}: नमक तथा
जल का दो-तिहाई, सारा ग्लूकोस तथा ऐमीनो अम्ल, और बाइकार्बोनेट पुनःसोखती
है। हेनले की अवरोही भुजा: मज्जा को जल खो देती है। आरोही भुजा: NaCl
पंप करती है, जल के लिए अपारगम्य। दूरस्थ नलिका: सोडियम साधती है
(\omterm{def:b3:renal-osmoregulation:controls}{ऐल्डोस्टेरॉन}), पोटैशियम स्रवित करती है। \omterm{def:b3:renal-osmoregulation:nephron}{संग्रह नलिका}: वैसोप्रेसिन के
तहत जल-पुनःसोखन, और प्रोटॉन तथा यूरिया की सँभाल।
\end{solution}

\begin{solution}{exo:b3:renal-osmoregulation:2}
निकासी: प्लाज़्मा का वह आयतन जो प्रति एकांक समय उस पदार्थ से साफ़ होता
है, $UV/P$। यहाँ $100\times 1.5/2 = \qty{75}{mL/min}$, जो
\qty{120}{mL/min} की GFR से नीचे है: वह पदार्थ छना जाता है और कुछ
पुनःसोखा जाता है (छनी मात्रा का लगभग \qty{38}{\%})।
\end{solution}

\begin{solution}{exo:b3:renal-osmoregulation:3}
कोई मीठे पानी की मछली अपने माध्यम से अधिक नमकीन है: जल क्लोमों से होकर
परासरण द्वारा भीतर आता है, इसलिए पीना उस बाढ़ में और जोड़ना भर होता; वह
भरपूर तनु मूत्र निकालती है और उसके क्लोमों की \omterm{def:b3:renal-osmoregulation:comparative}{क्लोराइड कोशिकाएँ} जल से
Na$^{+}$ तथा Cl$^{-}$ सक्रिय रूप से लेती हैं। कोई समुद्री मछली समुद्र
से कम नमकीन है: वह क्लोमों से जल खोती है, उसकी भरपाई के लिए समुद्री जल
पीती है, और उसकी \omterm{def:b3:renal-osmoregulation:comparative}{क्लोराइड कोशिकाएँ} निगला हुआ नमक बाहर पंप करती हैं।
\end{solution}

\begin{solution}{exo:b3:renal-osmoregulation:4}
\omterm{prop:b3:renal-osmoregulation:countercurrent}{एकल प्रभाव} वह \qty{200}{mOsm/L} का अंतर है जो मोटी आरोही भुजा अंतराल
और अपने ही तरल के बीच बनाए रखती है, NaCl बाहर पंप करके और जल को उसके
पीछे न आने देकर। आरोही भुजा को जल के लिए अपारगम्य होना ही पड़ता है,
वरना जल उस नमक के पीछे चला जाता और वह अंतर मिट जाता; अवरोही भुजा को जल
के लिए पारगम्य होना पड़ता है ताकि मोड़ तक पहुँचता तरल उस अंतराल से पहले
ही गाढ़ा हो चुका हो और हर चरण पिछले से नमकीन तरल पर काम करे। एक ही
पारगम्यता वाली दो भुजाएँ गुणा नहीं कर सकतीं।
\end{solution}

\begin{solution}{exo:b3:renal-osmoregulation:5}
अभिवाही सिरा: $60 - 18 - 28 = \qty{14}{mmHg}$; अपवाही सिरा: $60 - 18 -
35 = \qty{7}{mmHg}$। जैसे-जैसे जल छना जाता है, पीछे बचे प्रोटीन गाढ़े
होते हैं और उस केशिका के अनुदिश कलिलाकर्षी दाब चढ़ता है; वह
\qty{42}{mmHg} तक पहुँच जाए तो शुद्ध दाब शून्य हो जाता है और निस्यंदन
अंत से पहले ही रुक जाता है (निस्यंदन-संतुलन)। तब GFR प्लाज़्मा प्रवाह
पर निर्भर करती है: तेज़ प्रवाह प्रोटीनों को धीरे गाढ़ा करता है, इसलिए
निस्यंदन लंबे हिस्से तक चलता है और GFR प्रवाह के साथ चढ़ती है।
\end{solution}

\begin{solution}{exo:b3:renal-osmoregulation:6}
छना भार: $\qty{40}{mL/min}\times\qty{0.14}{mmol/mL} =
\qty{5.6}{mmol/min}$। उत्सर्जन: $70\times 1 = \qty{0.07}{mmol/min}$।
पुनःसोखा: $5.53/5.6 = \qty{98.75}{\%}$। सोडियम की निकासी: $70\times
1/140 = \qty{0.5}{mL/min}$।
\end{solution}

\begin{solution}{exo:b3:renal-osmoregulation:7}
दिन भर में $900/1200 = \qty{0.75}{L}$। \qty{400}{mOsm/L} पर: $900/400 =
\qty{2.25}{L}$ मूत्र, इसलिए लगभग $2.25 + 0.9 = \qty{3.15}{L}$ पीना
पड़ेगा (भोजन का जल घटाकर): जिस रोगी का गाढ़ा करने वाला तंत्र विफल हो
रहा हो उसे पीना ही पड़ता है, और न पी सके तो वह ख़तरे में है।
\end{solution}

\begin{solution}{exo:b3:renal-osmoregulation:8}
(क) डायबिटीज़ इन्सिपिडस: \qtyrange{50}{100}{mOsm/L} पर
\qtyrange{10}{15}{L} मूत्र, और जब भी पीना पिछड़े तब ऊँचा प्लाज़्मा
सोडियम, तथा न थमती प्यास। (ख) अनुचित वैसोप्रेसिन: थोड़ा गाढ़ा मूत्र,
जल रुका हुआ, तनुकरण से प्लाज़्मा सोडियम कम, और प्यास बढ़ी हुई नहीं
हालाँकि पीना चलता रहता है --- और सोडियम बहुत गिरे तो भ्रम तथा दौरे।
(ग) \omterm{def:b3:renal-osmoregulation:controls}{ऐल्डोस्टेरॉन} की अधिकता: सोडियम जल के साथ रुका, आयतन फैला, रक्तचाप
ऊँचा, पोटैशियम गया (कम प्लाज़्मा K$^{+}$), और प्लाज़्मा सोडियम लगभग
सामान्य, क्योंकि जल नमक के पीछे चलता है और गुर्दा कुछ बच भी निकलता है;
मूत्र-आयतन सामान्य। (घ) बिना नमक का आहार: रेनिन और \omterm{def:b3:renal-osmoregulation:controls}{ऐल्डोस्टेरॉन} दिन भर
में चढ़ जाते हैं, मूत्र का सोडियम दिन भर में कुछ मिलीमोल रह जाता है,
प्लाज़्मा सोडियम थमा रहता है, और आयतन तथा दाब थोड़े नीचे।
\end{solution}

\begin{solution}{exo:b3:renal-osmoregulation:9}
तनु मूत्र: परासरणी निकासी $100\times 6/300 = \qty{2}{mL/min}$;
\omterm{prop:b3:renal-osmoregulation:water}{मुक्त-जल निकासी} $6 - 2 = +\qty{4}{mL/min}$: गुर्दा विलेय-रहित जल बहा
रहा है, जैसा किसी जल-भार के बाद, जब वैसोप्रेसिन बंद हो।
गाढ़ा मूत्र: परासरणी निकासी $1000\times 0.6/300 =
\qty{2}{mL/min}$; \omterm{prop:b3:renal-osmoregulation:water}{मुक्त-जल निकासी} $0.6 - 2 = -\qty{1.4}{mL/min}$:
\qty{1.4}{mL/min} मुक्त जल शरीर को लौटाया जा रहा है --- यानी निर्जलन,
और ऊँचा वैसोप्रेसिन।
\end{solution}

\begin{solution}{exo:b3:renal-osmoregulation:10}
चरणों को वल्कुट से गिनिए; मान लीजिए चरण $k$ पर अंतराल
$I_{k}$ है। $k$ पर अवरोही तरल $I_{k}$ के बराबर है; $k$ पर आरोही तरल
$I_{k} - 200$ है; और $k$ पर आरोही तरल वही है जो मोड़ से चला और चरण
$n, n-1, \dots, k+1$ पार कर आया। स्थायी अवस्था में \qty{300}{} पर
घुसता तरल आरोही भुजा के शीर्ष से $I_{1} - 200$ पर निकलता है, और हर चरण
का अंतराल ऊपरवाले से अधिक से अधिक उस \omterm{prop:b3:renal-osmoregulation:countercurrent}{एकल प्रभाव} जितना आगे है, इसलिए
$I_{k} \le 300 + 200k$, और बराबरी तब जब हर चरण पूरे प्रभाव पर काम करे:
मोड़ $300 + 200n$ तक पहुँचता है। इसलिए वह प्रवणता पाश की लंबाई (यानी
चरणों की संख्या) गुणा \omterm{prop:b3:renal-osmoregulation:countercurrent}{एकल प्रभाव} से बँधी है, किसी कोशिका की
पंप-क्षमता से नहीं। $n$ पाश की लंबाई को उस दूरी से भाग देने पर मिलता है
जिस पर तरल अंतराल के साथ संतुलित होता है: लंबे उपमज्जीय पाश, और
रेगिस्तानी कृंतकों के बहुत लंबे पाश, बड़ा $n$ देते हैं; और पंप की
क्षमता तथा वासा रेक्टा द्वारा विलेय का बह जाना उस प्रभाव की छत बाँध
देते हैं।
\end{solution}

\begin{solution}{exo:b3:renal-osmoregulation:11}
\omterm{def:b3:renal-osmoregulation:comparative}{उपापचयी जल}: $5\times 0.8 = \qty{4}{g}$ स्टार्च $4\times 0.6
= \qty{2.4}{g}$ देता है। हानियाँ: मूत्र $3/5500\ \unit{L} = \qty{0.55}{mL}$;
वाष्पन \qty{1.5}{g} (जो नाक की वसूली के बाद का शुद्ध है)। कुल हानि
\qty{2.05}{g}, \qty{2.4}{g} के लाभ के सामने: यानी \qty{0.35}{g} का
अधिशेष, जो मल की छोटी हानि पूरी कर देता है, और बीजों का आर्द्रताग्राही
जल ऊपर से मिलता है। वह बिना पिए जी लेता है। नाक की वसूली के बिना
साँस की हानि $1.5/0.3 = \qty{5}{g}$ होती, और वह कुछ ही दिनों में मर
जाता।
\end{solution}

\begin{solution}{exo:b3:renal-osmoregulation:12}
समांतर प्रवाह में रक्त और अपोहक-द्रव उस झिल्ली के आधे रास्ते तक ही
संतुलित हो जाते हैं और विसरण चलाता सांद्रता-अंतर शून्य पर आ जाता है:
सबसे अच्छा हो तो रक्त अपनी आधी यूरिया के साथ निकलता है। प्रतिधारा
प्रवाह में सबसे ताज़ा अपोहक-द्रव सबसे साफ़ हुए रक्त से मिलता है और वह
प्रवणता पूरी झिल्ली भर बनी रहती है, इसलिए रक्त लगभग अपोहक-द्रव की
प्रवेश-सांद्रता के साथ निकल सकता है --- यानी वही सिद्धांत जो हेनले के
पाश का है। निकासी: $300\times 0.7 =
\qty{210}{mL/min}$, हफ़्ते में $3\times 4 = \qty{12}{h}$ के दौरान, यानी हफ़्ते में
$210\times 720 = \qty{151}{L}$; \qty{125}{mL/min} पर दो गुर्दे
$125\times\num{10080} = \qty{1260}{L}$ साफ़ करते हैं। वह मशीन गुर्दों
के काम का लगभग आठवाँ भाग करती है, यानी काल-औसत
\qty{15}{mL/min} --- जीने के लिए काफ़ी, स्वस्थ रहने के लिए नहीं।
\end{solution}

\begin{solution}{pb:b3:renal-osmoregulation:1}
\textbf{1.} शुद्ध दाब $55 - 15 - 30 = \qty{10}{mmHg}$; GFR $= 12.5
\times 10 = \qty{125}{mL/min}$।
\textbf{2.} $125\times 1440 = \qty{180}{L/d}$; $180/3 = 60$ बार।
\textbf{3.} इनुलिन: $62.5\times 1/0.5 = \qty{125}{mL/min}$, यानी GFR के
बराबर। क्रिएटिनिन: $1.3\times 1/0.01 = \qty{130}{mL/min}$, जो थोड़ी
बड़ी है क्योंकि \omterm{def:b3:renal-osmoregulation:nephron}{समीपस्थ नलिका} छने हुए के अलावा थोड़ा क्रिएटिनिन स्रवित
भी करती है।
\textbf{4.} PAH की निकासी $13\times 1/0.02 = \qty{650}{mL/min}$, यानी
वृक्कीय प्लाज़्मा प्रवाह; रक्त-प्रवाह $650/(1 - 0.45) = \qty{1180}{mL/min}$,
यानी लगभग \qty{1.2}{L/min}; निस्यंदन-अंश $125/650 = 0.19$।
\textbf{5.} शुद्ध दाब \qty{15}{mmHg}, GFR \qty{188}{mL/min} (असल में
कम, क्योंकि उस केशिका के अनुदिश कलिलाकर्षी दाब चढ़ता है)।
ऐंजियोटेंसिन~II अपवाही अणुधमनी सिकोड़ता है और केशिकागुच्छी दाब चढ़ा
देता है; उसे रोकने से वह दाब और GFR थोड़े गिर जाते हैं। मधुमेह में
\omterm{def:b3:renal-osmoregulation:nephron}{केशिकागुच्छ} ऊँचे दाब पर छानते हैं और वह छन्ना घिसता है; दाब गिराना उस
क्षति को वर्षों तक धीमा कर देता है।
\textbf{6.} पादकोशिकाएँ और उनके दरार-पट (या आधार कला का आवेश): यानी वह
परत जो ऐल्ब्युमिन रोकती है। \qty{3}{g/d} खोने से प्लाज़्मा का
कलिलाकर्षी दाब गिरता है; हर जगह केशिकाओं से तरल निकलता है और ऊतक सूज
जाते हैं --- यानी वृक्कोतक संलक्षण की शोफ।
\textbf{7.} छना $125\times 5 = \qty{0.625}{mmol/min}$; उत्सर्जित
$250\times 1 = \qty{0.25}{mmol/min}$; पुनःसोखा $0.375$, यानी
\qty{60}{\%}; निकासी $250/5 = \qty{50}{mL/min}$।
\textbf{8.} छना $180\times 0.14 = \qty{25.2}{mol/d}$ Na$^{+}$, यानी
\qty{580}{g} सोडियम, \qty{1.47}{kg} NaCl के रूप में। उत्सर्जित
$100\times 1.44
= \qty{144}{mmol/d}$; पुनःसोखा $1 - 144/\num{25200} = \qty{99.4}{\%}$।
दिन भर में लगभग डेढ़ किलोग्राम नमक खो जाता।
\textbf{9.} छना $125\times 5 = \qty{0.625}{mmol/min}$, यानी
\qty{112}{mg/min}, \qty{162}{g/d}; कुछ भी उत्सर्जित नहीं, क्योंकि यह
$T_{m}$ से नीचे है। छना भार $T_{m}$ तक $2.1/0.125 =
\qty{16.8}{mmol/L}$ (\qty{300}{mg/dL}) पर पहुँचता है --- जो लगभग
\qty{10}{mmol/L} की देखी गई देहली से ऊँचा है, क्योंकि \omterm{def:b3:renal-osmoregulation:nephron}{वृक्काणु}
एक-दूसरे से भिन्न हैं और सबसे कमज़ोर पहले छलकते हैं (यानी उस अनुमापन
वक्र का फैलाव)।
\textbf{10.} छना $125\times 20 = \qty{2.5}{mmol/min}$; उत्सर्जित
$2.5 - 2.1 = \qty{0.4}{mmol/min}$, यानी \qty{72}{mg/min}, दिन भर में
$576$ mmol या \qty{104}{g}। अतिरिक्त मूत्र $576/300 = \qty{1.9}{L}$।
वह ग्लूकोस जल को मूत्र में खींच ले जाता है (परासरणी मूत्रलता), रोगी
निर्जलित होता है और प्यासा; और \qty{104}{g} ग्लूकोस दिन भर में
\qty{1.8}{MJ} की हानि है, तथा वे कोशिकाएँ, जो ग्लूकोस ले नहीं पातीं,
वसा तथा प्रोटीन जलाती हैं: भार गिरता है।
\textbf{11.} पुनःसोखा $\approx \qty{25}{mol/d}$; ATP $25/3 =
\qty{8.4}{mol/d}$; $8.4\times 50 = \qty{420}{kJ}$, यानी \qty{7}{MJ} का
लगभग \qty{6}{\%}।
\textbf{12.} हर कचरा स्रवित करने के लिए उस नलिका को ऐसा कोई वाहक चाहिए
होता जो उसे पहचाने --- यानी हज़ारों वाहक, और किसी ऐसे अणु के लिए तो
कोई भी नहीं जो पहले कभी मिला ही न हो। सब छोटा-मोटा छान लेना और शरीर की
चाही कुछ सौ जातियाँ पुनःसोख लेना केवल उन्हीं के वाहक माँगता है जिनका
मूल्यवान होना पता है; और जो कुछ अपहचाना है वह अपने आप निकल जाता है।
कीमत प्रश्न~11 की ऊर्जा है।
\textbf{13.} परासरणी निकासी $600\times 1/290 = \qty{2.07}{mL/min}$;
\omterm{prop:b3:renal-osmoregulation:water}{मुक्त-जल निकासी} $1 - 2.07 = -\qty{1.07}{mL/min}$: गुर्दा जल बचा रहा है।
\textbf{14.} न्यूनतम $600/1200 = \qty{0.5}{L}$; अधिकतम $600/50 =
\qty{12}{L}$; न्यूनतम पर दिन भर की ज़रूरत $0.5 + 0.9 = \qty{1.4}{L}$।
\textbf{15.} \qty{1000}{mOsm} को \qty{1200}{mOsm/L} पर
\qty{0.83}{L} मूत्र चाहिए: उस दिन के विलेय से पहले शुद्ध लाभ
\qty{0.17}{L}। बाद में: उस दिन का कुल विलेय \qty{1600}{mOsm} है,
\qty{1.33}{L} मूत्र, जमा \qty{0.9}{L} अगोचर, यानी \qty{1}{L} भीतर के
लिए \qty{2.23}{L} बाहर --- यानी \qty{1.23}{L} की कमी, जबकि बिना पिए
\qty{1.4}{L} की होती। लीटर भर समुद्री जल उस संतुलन को केवल
\qty{0.17}{L} सुधारता है।
\textbf{16.} उस दिन के लिए \qty{50}{mOsm/L} पर मूत्र: $600/50 =
\qty{12}{L}$; रोज़ \qty{12.9}{L} पीना पड़ेगा, यानी लगभग \qty{13}{L}।
\textbf{17.} $B$ लीटर बियर $20B$ mOsm लाती है; अधिकतम मूत्र
$(600 + 20B)/50 = 12 + 0.4B$ लीटर है, और उसे $B - 0.9$
लीटर ढोना पड़ता है: $B - 0.9 \le 12 + 0.4B$, यानी $B \le \qty{21.5}{L}$।
लगभग बीस लीटर --- पर यदि और कुछ न खाया जाए तो उस दिन का विलेय गिरकर
\qty{200}{mOsm} रह जाता है और वह सीमा लगभग \qty{8}{L} पर आ जाती है:
यानी उन भारी पीने वालों की अल्पसोडियमता जो खाते नहीं।
\textbf{18.} $298/290 - 1 = \qty{2.8}{\%}$: यानी उस \qty{1}{\%} से कहीं
आगे जो वैसोप्रेसिन तथा प्यास छेड़ती है; दोनों अधिकतम हैं। अचर विलेय पर:
$290\times 42 = 298\times V$, $V = \qty{40.9}{L}$, यानी लगभग
\qty{1.1}{L} की कमी।
\textbf{19.} छह लीटर जल प्लाज़्मा तनु कर देता है जबकि पसीना नमक हटाता
है; और व्यायाम, पीड़ा तथा तनाव \omterm{def:b3:renal-osmoregulation:problem}{परासरणीयता} की परवाह किए बिना
वैसोप्रेसिन छोड़ते हैं, इसलिए गुर्दा उस अधिक जल को निकालने लायक तनु
मूत्र बना ही नहीं पाता। प्लाज़्मा सोडियम गिरता है, मस्तिष्क की कोशिकाएँ
फूलती हैं --- दौरे और मृत्यु तक हो चुके हैं। गुर्दे को वैसोप्रेसिन बंद
चाहिए होता और $U_{\min}$ पर \qty{6}{L} मूत्र; और उस धावक को प्यास के
हिसाब से पीना चाहिए तथा नमक की भरपाई करनी चाहिए।
\textbf{20.} दिन भर में $3/5500\ \unit{L} = \qty{0.55}{mL}$। किसी मानव
गुर्दे को उसी विलेय के लिए $3/1200 = \qty{2.5}{mL}$ चाहिए होते: कंगारू
चूहा उसका पाँचवाँ भाग जल काम में लेता है।
\textbf{21.} भीतर: \qty{2.4}{g}। बाहर: $1.6 + 0.3 + 0.55 = \qty{2.45}{g}$।
दिन भर में \qty{0.05}{g} के भीतर संतुलित, जिसे बीजों का आर्द्रताग्राही
जल (\qty{5}{g} का कुछ प्रतिशत) जुटा देता है।
\textbf{22.} जल $500\times 0.65 = \qty{325}{L}$; उसका चौथाई
\qty{81}{L}; \qty{5}{L/d} पर, सोलह दिन।
\textbf{23.} $150/800 = \qty{0.19}{L}$, यानी लगभग \qty{190}{mL} खारा
घोल --- \qty{100}{mL} समुद्री जल और उस मछली के जल से मिलकर कम, इसलिए वह
पक्षी जल पाता है। उसका गुर्दा, \qty{600}{mOsm/L} पर
(\qty{300}{mmol/L} NaCl), उसी नमक को ढोने के लिए \qty{0.5}{L} मूत्र
माँगता, यानी उससे अधिक जल जितना उसने लिया था: जो गुर्दा समुद्री जल को
हरा न सके वह किसी पक्षी को उसे पीने नहीं दे सकता, और वह ग्रंथि, जो
समुद्र से दुगुनी सांद्रता पर है, दे सकती है।
\textbf{24.} \qty{30}{g} जल भीतर, इसलिए दिन भर में लगभग \qty{30}{mL}
मूत्र, उसके रक्त की \omterm{def:b3:renal-osmoregulation:problem}{परासरणीयता} के किसी छोटे अंश पर (\num{300} के सामने
कुछ दसियों mOsm/L)। नमक उसी मूत्र में और क्लोमों के आर-पार विसरण से
निकलता है; और उस जल से, जिसमें प्रति लीटर कोई एक मिलीमोल है, \omterm{def:b3:renal-osmoregulation:comparative}{क्लोराइड
कोशिकाओं} के सक्रिय ग्रहण के बिना वह मछली दिनों में रिक्त हो जाती।
\textbf{25.} GFR \qty{125}{mL/min}, दिन भर में \qty{180}{L}; अनिवार्य
मूत्र दिन भर में \qty{0.5}{L}; और लीटर भर समुद्री जल, उसका नमक
उत्सर्जित हो जाने पर, शुद्ध \qty{0.17}{L} देता है --- यानी लगभग कुछ भी
नहीं।
\end{solution}
